\documentclass[conference]{IEEEtran}
\IEEEoverridecommandlockouts

%=========================== CODIFICACIÓN E IDIOMA ===========================
\usepackage[utf8]{inputenc}
\usepackage[T1]{fontenc}
\usepackage[spanish,es-tabla,es-nodecimaldot,es-noshorthands]{babel}
\usepackage{textcomp}

%=========================== MATEMÁTICAS Y UNIDADES ==========================
\usepackage{amsmath,amssymb,amsfonts}
\usepackage{siunitx}
\sisetup{output-decimal-marker={,}, group-separator={\,}, per-mode=symbol}
\DeclareSIUnit\dbi{dBi}
\DeclareSIUnit\sym{sym}
\DeclareSIUnit\samples{S}

%=========================== FIGURAS Y TABLAS ================================
\usepackage{graphicx}
\usepackage{float}
\usepackage{placeins}
\usepackage{stfloats}           % figuras de doble columna al pie [b]
\usepackage{needspace}          % mantiene juntos título y tabla del anexo
\renewcommand{\dblbotfraction}{0.75}
\usepackage{booktabs}
\usepackage{multirow}
\usepackage{caption}
\captionsetup{labelfont=bf, textfont=normalfont}
\usepackage{tikz}
\usetikzlibrary{arrows.meta,positioning,calc,fit,backgrounds}
\usepackage{pgfplots}
\pgfplotsset{compat=1.18}

%=========================== REFERENCIAS Y CÓDIGO ============================
\usepackage{url}
\usepackage{cite}
\usepackage{listings}
\usepackage{listingsutf8}
\usepackage{xcolor}
\definecolor{codebg}{RGB}{248,248,248}
\definecolor{codeframe}{RGB}{220,220,220}
\definecolor{codekw}{RGB}{0,92,197}
\definecolor{codestr}{RGB}{163,21,21}
\definecolor{codecmt}{RGB}{0,128,0}
\lstdefinestyle{bonito}{
    language=Python,
    backgroundcolor=\color{codebg},
    basicstyle=\ttfamily\small,
    keywordstyle=\color{codekw}\bfseries,
    stringstyle=\color{codestr},
    commentstyle=\color{codecmt}\itshape,
    numberstyle=\tiny\color{gray},
    numbers=left,
    stepnumber=1,
    frame=single,
    rulecolor=\color{codeframe},
    frameround=tttt,
    framesep=6pt,
    tabsize=2,
    showstringspaces=false,
    breaklines=true,
    upquote=true,
    xleftmargin=2em,
    framexleftmargin=2em,
    numbersep=8pt
}

%=========================== ESPACIADO =======================================
\raggedbottom
\usepackage{titlesec}
\titlespacing*{\section}{0pt}{1.2ex plus .3ex minus .2ex}{0.8ex}
\titlespacing*{\subsection}{0pt}{1.0ex plus .3ex minus .2ex}{0.6ex}
\titlespacing*{\subsubsection}{0pt}{0.8ex plus .2ex minus .2ex}{0.5ex}
\setlength{\parskip}{0pt}
\usepackage{enumitem}
\setlist{nosep}
\setlength{\abovecaptionskip}{4pt}
\setlength{\belowcaptionskip}{0pt}
\setlength{\textfloatsep}{10pt plus 1pt minus 2pt}
\setlength{\floatsep}{8pt plus 1pt minus 2pt}
\setlength{\intextsep}{8pt plus 1pt minus 2pt}
\renewcommand{\topfraction}{0.9}
\renewcommand{\bottomfraction}{0.8}
\renewcommand{\textfraction}{0.1}
\renewcommand{\floatpagefraction}{0.8}

%=========================== MACROS DEL PROYECTO =============================
% Texto pendiente (en rojo) para localizar fácilmente lo que falta completar.
\newcommand{\pendiente}[1]{\textcolor{red}{[#1]}}
% Nombre de las palabras clave en español (IEEEtran usa "Index Terms").
\renewcommand{\IEEEkeywordsname}{Palabras clave}
% Inserta una captura de la carpeta figures/. Si el archivo aún no existe,
% dibuja un recuadro rojo con el nombre esperado. Uso: \captura[ancho]{archivo.png}
\newcommand{\captura}[2][\linewidth]{%
  \IfFileExists{figures/#2}%
    {\includegraphics[width=#1]{figures/#2}}%
    {\fbox{\parbox[c][0.55\dimexpr#1\relax][c]{\dimexpr#1-2\fboxsep-2\fboxrule\relax}%
      {\centering\scriptsize\pendiente{\texttt{figures/\detokenize{#2}}}}}}}

\def\BibTeX{{\rm B\kern-.05em{\sc i\kern-.025em b}\kern-.08em
    T\kern-.1667em\lower.7ex\hbox{E}\kern-.125emX}}

\begin{document}

%%%%%%%%%%%%%%%%%%%%%%%%%%%%%%%%%%%%%%%%%%%%%%%%%%%%%%%%%%%%%%%%%%%%%%%%%%%%%
\title{Proyecto Integrador de Comunicaciones II:\\
Telemetría industrial bajo perturbaciones impulsivas}

\author{
\IEEEauthorblockN{Luis Miguel Higuera Ortiz}
\IEEEauthorblockA{\textit{Universidad Industrial de Santander}\\
Bucaramanga, Colombia\\
luis2230411@correo.uis.edu.co}
\and
\IEEEauthorblockN{Valeria Rueda Luna}
\IEEEauthorblockA{\textit{Universidad Industrial de Santander}\\
Bucaramanga, Colombia\\
Valeria2234069@correo.uis.edu.co}
\and
\IEEEauthorblockN{Maryuri Dayana Sanchez Sanchez}
\IEEEauthorblockA{\textit{Universidad Industrial de Santander}\\
Bucaramanga, Colombia\\
Maryuri2230398@correo.uis.edu.co}
\and
\IEEEauthorblockN{Jhoan Sebastian Vega Martinez}
\IEEEauthorblockA{\textit{Universidad Industrial de Santander}\\
Bucaramanga, Colombia\\
Jhoan2192280@correo.uis.edu.co}
}
%%%%%%%%%%%%%%%%%%%%%%%%%%%%%%%%%%%%%%%%%%%%%%%%%%%%%%%%%%%%%%%%%%%%%%%%%%%%%

\maketitle

% Cada sección se importa desde la carpeta sections/ para facilitar el trabajo
% conjunto. Los nombres de archivo no llevan tildes ni espacios (Overleaf).
%============================================================================
% RESUMEN
%============================================================================
\begin{abstract}
Se estudió un enlace de telemetría industrial de baja velocidad afectado por
perturbaciones impulsivas de motores, contactores y variadores. El sistema se
simuló en banda base con BPSK, pulso rectangular,
$R_s=\SI{32}{\kilo\sym\per\second}$, ocho muestras por símbolo y
$f_s=\SI{256}{\kilo\samples\per\second}$, con ruido gaussiano de fondo, un tren
periódico de impulsos de amplitud controlada y un receptor con filtro acoplado.
Se compararon una referencia y tres niveles de amplitud mediante la señal
temporal, la PSD, la constelación, el diagrama de ojo y la BER. No hubo errores
mientras la amplitud efectiva del impulso, diluida por el filtro acoplado, fue
menor que la del símbolo; por encima, la BER alcanzó un piso de
$1{,}1\times10^{-2}$, independiente de la amplitud y reproducido por el modelo
analítico. La PSD no mostró cambios apreciables, por lo que las medidas
promediadas no bastan para diagnosticar ruido impulsivo. Además, una prueba
de laboratorio con 16QAM entre dos radios HackRF One mostró que la constelación
solo fue reconocible al aumentar la relación señal a ruido acercando las placas.
\end{abstract}

\begin{IEEEkeywords}
BPSK, ruido impulsivo, telemetría industrial, filtro acoplado, tasa de error de bit.
\end{IEEEkeywords}


% INTRODUCCIÓN (presente/pasado, sin futuro)
%============================================================================
\section{Introducción}\label{sec:introduccion}
%============================================================================

El monitoreo de variables como temperatura, vibración o estado de máquina es
básico en la automatización industrial, y los enlaces inalámbricos que lo
soportan deben cumplir requisitos estrictos de confiabilidad~\cite{willig2005}
en entornos hostiles con interferencia de radiofrecuencia~\cite{gungor2009}. En
una planta, la conmutación de motores, contactores, variadores y cargas
inductivas produce interferencia impulsiva: eventos breves y de gran amplitud.
Los dispositivos con conmutación electromecánica, entre ellos los motores, se
han identificado como fuentes principales de este ruido~\cite{blackard1993}, que
es habitual en entornos industriales y degrada los sistemas
inalámbricos~\cite{chilo2009}; en canales sobre la red eléctrica incluye
componentes periódicas que hacen al canal variante en el tiempo~\cite{zimmermann2002}.

A diferencia del AWGN, estacionario y con potencia repartida en el tiempo, la
interferencia de origen humano es no gaussiana y degrada severamente a los
sistemas diseñados para ruido gaussiano~\cite{middleton1977}: concentra su
energía en pocas muestras de alta potencia~\cite{pighi2009}, y en portadora
única produce un piso de error que no desaparece al aumentar la energía de la
señal~\cite{ghosh1996}.

En este trabajo se estudió el escenario S1 del Proyecto Integrador de
Comunicaciones~II, representado en banda base: un enlace BPSK con pulso
rectangular, $R_s=\SI{32}{\kilo\sym\per\second}$, ocho muestras por símbolo y
$f_s=\SI{256}{\kilo\samples\per\second}$, al que se sumó un ruido impulsivo
periódico cuya amplitud fue la única variable, en una condición de referencia y
tres niveles de severidad.

El documento se organiza así: la Sección~\ref{sec:fundamento} presenta el
fundamento teórico; la Sección~\ref{sec:diseno}, el diseño del sistema; la
Sección~\ref{sec:metodologia}, la metodología de simulación; la
Sección~\ref{sec:resultados}, los resultados y su análisis; la
Sección~\ref{sec:hardware}, una prueba exploratoria en hardware SDR; y la
Sección~\ref{sec:conclusiones}, las conclusiones.

%----------------------------------------------------------------------------
\subsection{Objetivos}\label{subsec:objetivos}
%----------------------------------------------------------------------------

\textbf{General.} Evaluar en banda base la degradación de un enlace de
telemetría BPSK sometido a ruido impulsivo periódico de amplitud controlada.
\textbf{Específicos.} 1)~Justificar BPSK frente a OOK. 2)~Definir parámetros
coherentes ($R_s$, $R_b$, SPS, $f_s$ y ancho de banda). 3)~Modelar la
perturbación impulsiva y diferenciarla del AWGN. 4)~Comparar una referencia y
tres niveles de amplitud mediante señal temporal, PSD, constelación, diagrama
de ojo y BER, y contrastarlos con la hipótesis y el modelo analítico.


% FUNDAMENTO TEÓRICO, DISEÑO DEL SISTEMA Y METODOLOGÍA
%============================================================================
\section{Fundamento teórico}\label{sec:fundamento}
%============================================================================

\subsection{Modulación BPSK}\label{subsec:bpsk}

En BPSK cada bit se asocia a uno de dos símbolos antipodales. En banda base, la
señal con pulso de conformación $p(t)$ es
\begin{equation}
  s(t)=A\sum_{k} a_k\,p(t-kT_s), \qquad a_k\in\{-1,+1\},
  \label{eq:bpsk}
\end{equation}
donde $A$ es la amplitud y $T_s=1/R_s$ el periodo de símbolo. La información
está solo en la componente en fase (I); con $M=2$, $k=\log_2M=1$ y $R_b=R_s$.
Con pulso rectangular, $E_b=A^2T_s$, los símbolos están separados $2\sqrt{E_b}$
y, en AWGN con filtro acoplado y umbral en cero~\cite{proakis2008},
\begin{equation}
  P_b^{\mathrm{AWGN}}=Q\!\left(\sqrt{2E_b/N_0}\right).
  \label{eq:pb_awgn}
\end{equation}
Su PSD es $S_s(f)=A^2T_s\,\mathrm{sinc}^2(fT_s)$, con nulos en $f=\pm mR_s$, de
modo que el ancho de banda entre primeros nulos es $2R_s$~\cite{sklar2001}. El
lóbulo principal contiene cerca del \SI{90}{\percent} de la potencia y el primer
lóbulo lateral está unos \SI{13,3}{\decibel} por debajo; con un pulso RRC de
factor $\beta$, el ancho de banda sería $(1+\beta)R_s$.

\textbf{Alternativa descartada.} En OOK, con igual energía por bit, los símbolos
$0$ y $\sqrt{2E_b}$ están un factor $\sqrt2$ más cerca
($P_b=Q(\sqrt{E_b/N_0})$, \SI{3}{\decibel} de penalización)~\cite{proakis2008}, y
su umbral depende del nivel recibido, por lo que un impulso lo cruza con más
facilidad. Por ello se eligió BPSK.

\subsection{Ruido impulsivo frente a AWGN}\label{subsec:ruido}

El AWGN es gaussiano, estacionario y afecta por igual a todas las muestras. La
interferencia de origen humano es no gaussiana: los modelos de Middleton la
describen como una mezcla dominada por pocos eventos intensos~\cite{middleton1977},
y en tiempo discreto se usa el modelo Bernoulli-gaussiano~\cite{ghosh1996}. En
este trabajo se empleó un tren periódico controlado de impulsos rectangulares
de amplitud $A_i$, ancho $W$ y periodo $N_T$ muestras,
\begin{equation}
  i[n]=A_i\sum_{k}\Pi_W[n-kN_T],
  \label{eq:tren}
\end{equation}
donde $\Pi_W[n]=1$ para $0\le n<W$. Representa las componentes periódicas del
ruido de conmutación~\cite{zimmermann2002} y permite variar solo la amplitud.
Su ciclo útil es $\delta=W/N_T$ y su potencia media $A_i^2\delta$: con $\delta$
pequeño, la energía se concentra en pocas muestras y se reparte en una banda
ancha, con envolvente espectral de primer nulo en $f_s/W$ y líneas espaciadas
$f_T=f_s/N_T$; al ser unipolar, el tren también tiene una componente de continua
$A_i\delta$. La Tabla~\ref{tab:awgn_vs_imp} resume las diferencias con el AWGN.

\begin{table}[!t]
\centering
\caption{Comparación entre AWGN y ruido impulsivo periódico}
\label{tab:awgn_vs_imp}
\footnotesize
\setlength{\tabcolsep}{3pt}
\begin{tabular}{@{}p{0.24\columnwidth}p{0.35\columnwidth}p{0.35\columnwidth}@{}}
\toprule
\textbf{Aspecto} & \textbf{AWGN} & \textbf{Impulsivo periódico} \\
\midrule
Distribución temporal & Continua, todas las muestras & Concentrada en $W$ de cada $N_T$ muestras \\
Amplitud & Gaussiana, media cero & Pico $A_i$, alta relación pico/media \\
PSD & Piso plano & Líneas cada $f_T$ con envolvente ancha \\
Constelación & Nube circular en ambos símbolos & Puntos aislados desplazados $A_i$ \\
Error en BPSK & Gradual con $E_b/N_0$ & Umbral en $A_i W/\mathrm{SPS}\approx A$; piso de error \\
\bottomrule
\end{tabular}
\end{table}

\begin{figure}[!t]
\centering
\captura[0.95\columnwidth]{escenario.png}
\caption{Escenario S1: nodo Tx sobre la máquina, estación Rx de supervisión y
fuentes de perturbación impulsiva (motores, contactores y variadores).
Ilustración elaborada con apoyo de IA (Claude, Anthropic), que generó el código
de dibujo; el contenido técnico fue revisado por el equipo.}
\label{fig:escenario}
\end{figure}

\begin{figure*}[!t]
\centering
\begin{tikzpicture}[
  font=\scriptsize,
  >={Stealth[length=4pt]},
  blk/.style={draw, rounded corners=2pt, align=center, minimum height=1.15cm,
              inner sep=2pt, fill=blue!4, text width=1.6cm},
  rx/.style={blk, fill=orange!8},
  ch/.style={blk, fill=red!6, minimum height=0.85cm, text width=1.9cm},
  sum/.style={draw, circle, inner sep=0pt, minimum size=0.45cm, fill=white},
  obs/.style={draw=green!50!black, fill=green!8, rounded corners=1pt, font=\tiny,
              align=center, inner sep=1.5pt},
  tt/.style={font=\ttfamily\tiny}
]
  % Transmisor
  \node[blk] (src) {Fuente binaria\\[1pt] \texttt{\tiny Random Source}\\[-1pt]\texttt{\tiny Pack/Unpack}};
  \node[blk, right=0.35cm of src] (map) {Mapeo BPSK\\ $\{-1,+1\}$\\[1pt] \texttt{\tiny Chunks to Symbols}};
  \node[blk, right=0.35cm of map] (int) {Pulso rectangular\\ SPS $=8$\\[1pt] \texttt{\tiny Interp. FIR Filter}};
  \node[sum, right=1.0cm of int] (add) {$+$};
  % Canal
  \node[ch, above=0.45cm of add] (awgn) {AWGN de fondo\\ \texttt{\tiny noise\_amp} $=0{,}10$};
  \node[ch, below=0.45cm of add] (imp) {Tren de impulsos\\ $A_i$, $W=4$, $N_T=500$};
  % Receptor
  \node[rx, right=2.3cm of add] (mf) {Filtro acoplado\\ promedio de 8\\[1pt] \texttt{\tiny Moving Average}};
  \node[rx, right=0.35cm of mf] (smp) {Muestreo\\ 1 de cada 8\\[1pt] \texttt{\tiny Keep One in N}};
  \node[rx, right=0.35cm of smp] (dec) {Decisión\\ $\mathrm{Re}\{z_k\}\gtrless 0$\\[1pt] \texttt{\tiny Binary Slicer}};
  \node[rx, right=0.35cm of dec, text width=1.45cm] (ber) {BER\\[1pt] \texttt{\tiny XOR +}\\[-1pt]\texttt{\tiny Moving Average}};

  \draw[->] (src) -- node[above=1pt]{$R_b$} (map);
  \draw[->] (map) -- node[above=1pt]{$R_s$} (int);
  \draw[->] (int) -- node[above=1pt]{$f_s$} (add);
  \draw[->] (awgn) -- (add);
  \draw[->] (imp) -- (add);
  \draw[->] (add) -- (mf);
  \draw[->] (mf) -- (smp);
  \draw[->] (smp) -- (dec);
  \draw[->] (dec) -- (ber);

  % Bits de referencia
  \coordinate (yref) at ([yshift=-0.35cm]imp.south);
  \draw[->, dashed, gray] (src.south) -- (src.south |- yref) -- (ber.south |- yref) -- (ber.south);
  \node[above=0pt, text=black!70, font=\tiny] at ($(src.south |- yref)!0.5!(int.south |- yref)$)
       {bits transmitidos (referencia para la BER)};

  % Bloques del canal
  \begin{scope}[on background layer]
    \node[draw=red!50, dashed, rounded corners=3pt, inner sep=4pt, fit=(awgn)(imp)(add),
          label={[text=red!60!black, font=\tiny, anchor=south]north:Canal (banda base)}] {};
    \node[draw=blue!40, dashed, rounded corners=3pt, inner sep=4pt, fit=(src)(map)(int),
          label={[text=blue!60!black, font=\tiny, anchor=south]north:Transmisor}] {};
    \node[draw=orange!60, dashed, rounded corners=3pt, inner sep=4pt, fit=(mf)(smp)(dec)(ber),
          label={[text=orange!60!black, font=\tiny, anchor=south]north:Receptor}] {};
  \end{scope}

  % Puntos de observacion
  \node[obs, below=0.12cm of mf] (o2) {Ojo};
  \node[obs, below=0.12cm of smp] (o3) {Constelación\\ de decisión};
  \node[obs] at ($(add.east)!0.62!(mf.west)+(0,0.75)$) (o1) {Tiempo, PSD,\\ constelación};
  \draw[green!50!black, densely dotted] (o1.south) -- ($(add.east)!0.62!(mf.west)$);
\end{tikzpicture}
\caption{Diagrama técnico del sistema Tx--canal--Rx en banda base:
$R_s=\SI{32}{\kilo\sym\per\second}$, $R_b=\SI{32}{\kilo\bit\per\second}$,
SPS $=8$ y $f_s=\SI{256}{\kilo\samples\per\second}$. Los recuadros verdes
indican los puntos de observación. Implementación:
\texttt{S1\_BPSK\_Ruido\_Impulsivo\_Grupo3.grc} (bloques agregados en la Tabla~\ref{tab:bloques}).}
\label{fig:bloques}
\end{figure*}

\subsection{Efecto del impulso sobre la decisión BPSK}\label{subsec:efecto}

Con pulso rectangular, el filtro acoplado promedia las SPS muestras de cada
símbolo. La variable de decisión del símbolo $k$ es
\begin{equation}
  z_k=\frac{1}{\mathrm{SPS}}\sum_{n\in k} r[n] = A\,a_k + A_{\mathrm{ef}} + \tilde{w}_k,
  \label{eq:zk}
\end{equation}
donde $\tilde{w}_k$ tiene desviación $\sigma_z=\sigma_w/\sqrt{\mathrm{SPS}}$ y
$A_{\mathrm{ef}}=A_iW/\mathrm{SPS}$ es la amplitud efectiva de un impulso que cae
completo en el símbolo: el filtro acoplado lo diluye. Si una fracción
$\rho=\mathrm{SPS}/N_T$ de los símbolos contiene un impulso positivo y $p_-$ es la
fracción de ellos que vale $-1$ (1/2 con bits equiprobables),
\begin{equation}
  P_b\approx(1-\rho)\,Q\!\left(\tfrac{A}{\sigma_z}\right)
  +\rho\left[p_-Q\!\left(\tfrac{A-A_{\mathrm{ef}}}{\sigma_z}\right)
  +(1-p_-)Q\!\left(\tfrac{A+A_{\mathrm{ef}}}{\sigma_z}\right)\right].
  \label{eq:pb_imp}
\end{equation}
Con $\sigma_z\ll A$, la BER es nula si $A_i<A\,\mathrm{SPS}/W$, vale
$\rho p_-/2$ en el umbral y alcanza el piso $\rho p_-$ por encima de él:
\begin{equation}
  P_b\approx 0 \;\;(A_i<A\,\mathrm{SPS}/W),\qquad
  P_b\approx\rho\,p_- \;\;(A_i>A\,\mathrm{SPS}/W).
  \label{eq:pb_limite}
\end{equation}
El comportamiento es de umbral y el piso no depende de la energía de la señal,
como se ha reportado para sistemas de portadora única~\cite{ghosh1996}. Como
comparación, un receptor que tomara una sola muestra por símbolo, sin filtro
acoplado, tendría el umbral en $A_i=A$, pero solo fallaría cuando el impulso
coincide con el instante de muestreo, con $\rho\approx W/N_T$. El filtro acoplado
eleva entonces el umbral de amplitud a cambio de que todo impulso influya en
alguna decisión. Si el impulso tuviera componente en cuadratura, dispersaría la
constelación en vertical sin alterar la decisión BPSK, que depende solo de la
parte real: la dispersión visible no equivale necesariamente a errores.

%============================================================================
\section{Diseño del sistema y parámetros}\label{sec:diseno}
%============================================================================

\subsection{Escenario y parámetros}\label{subsec:escenario}

Un nodo sobre una máquina transmite variables de baja velocidad a la estación
de supervisión, en una nave con motores, contactores y variadores cuyas
conmutaciones generan perturbaciones breves de gran amplitud
(Fig.~\ref{fig:escenario}).

Se verificó que el paquete de referencia del curso fuera adecuado: una
vibración muestreada a \SI{1}{\kilo\samples\per\second} con 16 bits requiere
\SI{16}{\kilo\bit\per\second}, la mitad de la tasa disponible, y la temperatura
o el estado requieren mucho menos. Las relaciones de diseño son
\begin{align}
  k &= \log_2 2 = 1, \qquad R_b = R_s\,k = \SI{32}{\kilo\bit\per\second}, \label{eq:rb}\\
  f_s &= R_s\cdot\mathrm{SPS} = \num{32000}\times 8 = \SI{256}{\kilo\samples\per\second}, \label{eq:fs}\\
  B_{\mathrm{nulos}} &\approx 2R_s=\SI{64}{\kilo\hertz}, \label{eq:bw}
\end{align}
coherentes con el flujograma base. La ventana de la PSD
($\pm f_s/2=\pm\SI{128}{\kilo\hertz}$) contiene el lóbulo principal y tres
lóbulos laterales por lado. La Tabla~\ref{tab:parametros} resume el diseño.

\begin{table}[!t]
\centering
\caption{Parámetros de diseño del sistema}
\label{tab:parametros}
\footnotesize
\begin{tabular}{@{}llr@{}}
\toprule
\textbf{Parámetro} & \textbf{Símbolo} & \textbf{Valor} \\
\midrule
Modulación              & --                     & BPSK ($M=2$) \\
Bits por símbolo        & $k$                    & 1 \\
Tasa de símbolos        & $R_s$                  & \SI{32}{\kilo\sym\per\second} \\
Tasa de bits            & $R_b$                  & \SI{32}{\kilo\bit\per\second} \\
Periodo de símbolo      & $T_s=T_b$              & \SI{31,25}{\micro\second} \\
Muestras por símbolo    & SPS                    & 8 \\
Frecuencia de muestreo  & $f_s$                  & \SI{256}{\kilo\samples\per\second} \\
Conformación de pulso   & $p(t)$                 & Rectangular \\
Ancho de banda (nulos)  & $B_{\mathrm{nulos}}$   & \SI{64}{\kilo\hertz} \\
Ancho de banda RRC (ref.)& $B_{\mathrm{RRC}}$    & \SI{43,2}{\kilo\hertz} \\
Eficiencia espectral    & $R_b/B_{\mathrm{nulos}}$ & \SI{0,5}{\bit\per\second\per\hertz} \\
Perturbación            & $i[n]$                 & Impulsiva periódica \\
Variable experimental   & $A_i$                  & Amplitud del impulso \\
\bottomrule
\end{tabular}
\end{table}

\subsection{Arquitectura del sistema}\label{subsec:arquitectura}

En la Fig.~\ref{fig:bloques}, la fuente binaria alimenta el mapeo BPSK y el
interpolador de pulso rectangular; el canal suma un AWGN de fondo fijo y el
tren de~\eqref{eq:tren}; el receptor aplica el filtro acoplado, toma una
muestra por símbolo, decide por signo y compara con los bits transmitidos
para obtener la BER.

%============================================================================
\section{Metodología de simulación}\label{sec:metodologia}
%============================================================================

\begin{figure}[!t]
\centering
\captura[\columnwidth]{grc_base.png}
\caption{Flujograma base del curso con fuente de AWGN.}
\label{fig:grc_base}
\end{figure}

\begin{figure}[!t]
\centering
\captura[\columnwidth]{grc_modificado.png}
\caption{Flujograma modificado con generador de impulsos y receptor BPSK.}
\label{fig:grc_mod}
\end{figure}

\begin{figure*}[!t]
\centering
\setlength{\tabcolsep}{1pt}
\begin{tabular}{cccc}
\captura[0.24\textwidth]{tiempo_A.png} & \captura[0.24\textwidth]{tiempo_B.png} &
\captura[0.24\textwidth]{tiempo_C.png} & \captura[0.24\textwidth]{tiempo_D.png} \\
\scriptsize (a) A: $A_i=0$ & \scriptsize (b) B: $A_i=1$ &
\scriptsize (c) C: $A_i=2$ & \scriptsize (d) D: $A_i=4$ \\[2pt]
\captura[0.24\textwidth]{psd_A.png} & \captura[0.24\textwidth]{psd_B.png} &
\captura[0.24\textwidth]{psd_C.png} & \captura[0.24\textwidth]{psd_D.png} \\
\scriptsize (e) A: $A_i=0$ & \scriptsize (f) B: $A_i=1$ &
\scriptsize (g) C: $A_i=2$ & \scriptsize (h) D: $A_i=4$ \\
\end{tabular}
\caption{Señal en el canal (componentes I y Q; amplitud frente a tiempo en ms,
ventana de \SI{8}{\milli\second}; la escala vertical se ajustó en cada prueba)
en (a)--(d), y PSD (dB frente a frecuencia en kHz, $\pm\SI{128}{\kilo\hertz}$)
en (e)--(h), para las cuatro pruebas.}
\label{fig:tiempo_psd}
\end{figure*}

\begin{figure*}[!t]
\centering
\setlength{\tabcolsep}{1pt}
\begin{tabular}{cccc}
\captura[0.24\textwidth]{ojo_A.png} & \captura[0.24\textwidth]{ojo_B.png} &
\captura[0.24\textwidth]{ojo_C.png} & \captura[0.24\textwidth]{ojo_D.png} \\
\scriptsize (a) A: $A_i=0$ & \scriptsize (b) B: $A_i=1$ &
\scriptsize (c) C: $A_i=2$ & \scriptsize (d) D: $A_i=4$ \\[2pt]
\captura[0.24\textwidth]{const_A.png} & \captura[0.24\textwidth]{const_B.png} &
\captura[0.24\textwidth]{const_C.png} & \captura[0.24\textwidth]{const_D.png} \\
\scriptsize (e) A: $A_i=0$ & \scriptsize (f) B: $A_i=1$ &
\scriptsize (g) C: $A_i=2$ & \scriptsize (h) D: $A_i=4$ \\
\end{tabular}
\caption{Diagrama de ojo a la salida del filtro acoplado (componente I; amplitud
frente a tiempo en ms) en (a)--(d), y constelación en el instante de decisión
(I frente a Q) en (e)--(h), para las cuatro pruebas.}
\label{fig:ojo_const}

\vspace{6pt}
\setlength{\tabcolsep}{1pt}
\begin{tabular}{cccc}
\captura[0.24\textwidth]{const_sobre_A.png} & \captura[0.24\textwidth]{const_sobre_B.png} &
\captura[0.24\textwidth]{const_sobre_C.png} & \captura[0.24\textwidth]{const_sobre_D.png} \\
\scriptsize (a) A: $A_i=0$ & \scriptsize (b) B: $A_i=1$ &
\scriptsize (c) C: $A_i=2$ & \scriptsize (d) D: $A_i=4$ \\
\end{tabular}
\caption{Constelación sobremuestreada a la salida del canal (todas las
muestras, I frente a Q) para las cuatro pruebas.}
\label{fig:const_sobre}
\end{figure*}

\subsection{Flujograma base y modificaciones}\label{subsec:grc_base}

Se partió de \texttt{Modulaciones\_basadas\_en\_constelaciones\_AWGN.grc}, que
conserva el transmisor del archivo base y suma ruido gaussiano
(Fig.~\ref{fig:grc_base}). Su fuente aleatoria repite una secuencia de 1000
bits, el bloque \texttt{Chunks to Symbols} usa la tabla de constelación y el
interpolador con \texttt{h=[1]*Sps} forma el pulso rectangular. El original se
conservó y se trabajó sobre la copia
\texttt{S1\_BPSK\_Ruido\_Impulsivo\_Grupo3.grc}. Como la tabla por defecto era
la de OOK, $(0,1)$, se cambió por la de BPSK, $(-1,+1)$, y se verificó que
$M=2$ y $R_b=\SI{32}{\kilo\bit\per\second}$.

La Tabla~\ref{tab:bloques} y la Fig.~\ref{fig:grc_mod} muestran los bloques
agregados. El tren de impulsos se generó con un \texttt{Vector Source} de
$N_T$ muestras, $W$ de ellas en uno, multiplicado por \texttt{A\_imp} y sumado
en una tercera entrada del \texttt{Add}. El receptor usa un promedio móvil de
SPS muestras como filtro acoplado, un \texttt{Keep One in N} que conserva la
última muestra de cada símbolo (máximo del integrador) y un
\texttt{Binary Slicer}, que entrega 1 si la parte real es positiva, de acuerdo
con la asignación $0\to-1$ y $1\to+1$. La BER se obtiene con una XOR entre bits
transmitidos y decididos, cuya salida vale 1 en cada error, promediada sobre
$10^5$ bits (unos \SI{3,1}{\second}). Como ni el interpolador ni el promedio
móvil introducen retardo de símbolo, ambas secuencias quedan alineadas; la
alineación se validó con BER nula en la referencia.

\begin{table}[!t]
\centering
\caption{Bloques agregados al flujograma}
\label{tab:bloques}
\scriptsize
\setlength{\tabcolsep}{3pt}
\begin{tabular}{@{}p{0.30\columnwidth}p{0.30\columnwidth}p{0.34\columnwidth}@{}}
\toprule
\textbf{Bloque} & \textbf{Parámetros} & \textbf{Función} \\
\midrule
\texttt{Vector Source} (complex) & \texttt{[1]*W\_imp + [0]*(N\_T-W\_imp)}, repetir & Forma del tren de impulsos \\
\texttt{Multiply Const} & \texttt{A\_imp} & Amplitud del impulso (variable) \\
\texttt{Add} (3 entradas) & -- & Señal + AWGN + impulsos \\
\texttt{Moving Average} (complex) & longitud 8, escala $1/8$ & Filtro acoplado \\
\texttt{QT GUI Eye Sink} & 8 muestras/símbolo & Diagrama de ojo \\
\texttt{Keep One in N} & $N=8$ & Una muestra por símbolo \\
\texttt{QT GUI Constellation Sink} & -- & Constelación de decisión \\
\texttt{Complex to Real} + \texttt{Binary Slicer} & umbral 0 & Decisión de bit \\
\texttt{XOR} + \texttt{Char to Float} & -- & Detección de errores \\
\texttt{Moving Average} (float) & longitud $10^5$ & Estimación de BER \\
\texttt{QT GUI Number Sink} & -- & Visualización de BER \\
\bottomrule
\end{tabular}
\end{table}

\subsection{Plan de pruebas}\label{subsec:plan}

Se fijaron $W=4$ muestras (\SI{15,6}{\micro\second}, medio símbolo) y
$N_T=500$ (\SI{1,953}{\milli\second}). La repetición es más rápida que la de una
conmutación real para observar cuatro impulsos por captura y reunir errores en
pocos segundos. Como 500 no es múltiplo de 8, los impulsos caen alternadamente
en la primera y la segunda mitad de un símbolo, pero siempre completos dentro
de él, de modo que
$\rho=\mathrm{SPS}/N_T=0{,}016$ y $A_{\mathrm{ef}}=A_i/2$. El AWGN de fondo
(\texttt{noise\_amp} $=0{,}10$, $\sigma_z=0{,}025$) no produce errores por sí
solo. Solo cambió $A_i$, con niveles alrededor del umbral teórico
$A\,\mathrm{SPS}/W=2$ (Tabla~\ref{tab:pruebas}, con
$\mathrm{SIR}=A^2/(A_i^2\delta)$ y la BER esperada para bits equiprobables). En cada prueba se esperaron
al menos \SI{5}{\second} antes de registrar las gráficas y la BER.

\begin{table}[!t]
\centering
\caption{Plan de pruebas: condición de referencia y tres niveles de perturbación}
\label{tab:pruebas}
\footnotesize
\setlength{\tabcolsep}{4pt}
\begin{tabular}{@{}clcccc@{}}
\toprule
\textbf{Prueba} & \textbf{Severidad} & \textbf{$A_i$} & \textbf{$A_{\mathrm{ef}}$} &
\textbf{SIR [dB]} & \textbf{BER esperada} \\
\midrule
A & Referencia & 0   & 0   & $\infty$ & 0 \\
B & Suave      & 1,0 & 0,5 & 21,0 & $\approx 0$ \\
C & Media      & 2,0 & 1,0 & 14,9 & $\approx 4\times10^{-3}$ \\
D & Fuerte     & 4,0 & 2,0 & 8,9  & $\approx 8\times10^{-3}$ \\
\midrule
\multicolumn{6}{@{}p{0.95\columnwidth}@{}}{Constantes: $W=4$, $N_T=500$,
\texttt{noise\_amp} $=0{,}10$, $N_{\mathrm{ber}}=10^5$; $R_s$, SPS y $f_s$ según la
Tabla~\ref{tab:parametros}.} \\
\bottomrule
\end{tabular}
\end{table}

\subsection{Hipótesis}\label{subsec:hipotesis}

Antes de simular se esperaba que mayores amplitudes produjeran perturbaciones
temporales más fuertes y mayor dispersión. Con el modelo anterior: impulsos
visibles desde B sin errores; en C, BER cercana a $\rho/4=4\times10^{-3}$; en D,
un piso de $\rho/2=8\times10^{-3}$ independiente de la potencia de la señal; y
en la PSD, un nivel adicional de banda ancha creciente con $A_i^2$.


% RESULTADOS Y ANÁLISIS
%============================================================================
\section{Resultados y análisis}\label{sec:resultados}
%============================================================================

\subsection{Validación del enlace (prueba A)}\label{subsec:validacion}

En la referencia, la componente I alternó entre $\pm1$ con el leve
ensanchamiento del AWGN y Q contuvo solo ruido (Fig.~\ref{fig:tiempo_psd}a). La
PSD mostró el lóbulo principal centrado en $f=0$ y lóbulos laterales
decrecientes (Fig.~\ref{fig:tiempo_psd}e). El ojo presentó rampas lineales,
propias del integrador, con vértices en $\pm1$ cada \SI{31,25}{\micro\second}
(Fig.~\ref{fig:ojo_const}a), es decir, una apertura vertical de 2 en el
instante de decisión. Las nubes de la constelación sobremuestreada
($\approx\pm0{,}2$, Fig.~\ref{fig:const_sobre}a) se redujeron a
$\approx\pm0{,}07$ en la de decisión (Fig.~\ref{fig:ojo_const}e): el factor
$\sqrt8$ es la ganancia del filtro acoplado. La BER fue cero, lo que validó el
receptor y la alineación de los bits.

\subsection{Señal temporal y PSD}\label{subsec:res_tiempo_psd}

En D aparecieron cuatro picos, cerca de \num{1,5}, \num{3,45}, \num{5,4} y
\SI{7,35}{\milli\second}, separados unos \SI{1,95}{\milli\second} como en el
diseño ($N_T/f_s=\SI{1,953}{\milli\second}$), con alturas $1+A_i=5$ sobre símbolos $+1$ y $-1+A_i=3$ sobre $-1$
(Fig.~\ref{fig:tiempo_psd}d). En B y C solo fueron evidentes los impulsos sobre
símbolos $+1$; sobre un $-1$ llevaron la señal a 0 o a $+1$ durante
\SI{15,6}{\micro\second} y se confundieron con una transición normal: una
perturbación de este nivel puede imitar un símbolo válido. Q no se alteró,
porque el impulso es real. Las cuatro PSD conservaron la misma forma, con
diferencias del orden de las fluctuaciones de cada captura
(Fig.~\ref{fig:tiempo_psd}e--h). Con $\delta=0{,}008$ la potencia media del
impulso es baja y se reparte en una banda ancha, de modo que la PSD, una medida
promediada, resultó poco sensible a esta perturbación escasa e intensa, a diferencia del
AWGN, que eleva el piso espectral de forma visible~\cite{middleton1977,pighi2009}.

\subsection{Constelación y diagrama de ojo}\label{subsec:res_const_ojo}

En la constelación sobremuestreada aparecieron pocos puntos desplazados
exactamente $A_i$: en 0 y 2 (B), 1 y 3 (C), y 3 y 5 (D)
(Fig.~\ref{fig:const_sobre}b--d), pues solo 4 de cada 500 muestras están
perturbadas. En la de decisión el desplazamiento fue la mitad: $-0{,}5$ y
$1{,}5$ (B), 0 y 2 (C), y $+1$ y 3 (D) (Fig.~\ref{fig:ojo_const}f--h), lo que
confirma $A_{\mathrm{ef}}=A_i/2$. En C, los símbolos $-1$ alcanzados quedaron
sobre el umbral; en D se superpusieron a la nube de $+1$ y se decidieron mal. En
el ojo, la mayoría de los trazos no cambió y el margen de los $-1$ alcanzados,
$1-A_{\mathrm{ef}}$, pasó de 1 a $0{,}5$, 0 y $-1$ en B, C y D
(Fig.~\ref{fig:ojo_const}b--d): el cierre no fue uniforme, como con AWGN, sino
limitado al $1{,}6\,\%$ de símbolos afectados.

\subsection{Tasa de error de bit}\label{subsec:res_ber}

\begin{table}[!t]
\centering
\caption{BER medida frente a la BER teórica de~\eqref{eq:pb_imp}, con bits
equiprobables ($p_-=1/2$) y con la secuencia periódica empleada ($p_-=11/16$)}
\label{tab:ber}
\footnotesize
\setlength{\tabcolsep}{3.5pt}
\begin{tabular}{@{}cccccc@{}}
\toprule
\textbf{Prueba} & \textbf{$A_i$} & \textbf{$1-A_{\mathrm{ef}}$} &
\textbf{$p_-=1/2$} & \textbf{$p_-=11/16$} & \textbf{Medida} \\
\midrule
A & 0   & 1   & 0                    & 0                    & 0 \\
B & 1,0 & 0,5 & $\approx 0$          & $\approx 0$          & 0 \\
--& 1,8 & 0,1 & $2{,}5\times10^{-7}$ & $3{,}5\times10^{-7}$ & 0 \\
--& 1,9 & 0,05& $1{,}8\times10^{-4}$ & $2{,}5\times10^{-4}$ & $2{,}8\times10^{-4}$ \\
C & 2,0 & 0   & $4{,}0\times10^{-3}$ & $5{,}5\times10^{-3}$ & $5{,}82\times10^{-3}$ \\
--& 2,1 & $-0{,}05$ & $7{,}8\times10^{-3}$ & $1{,}07\times10^{-2}$ & $1{,}07\times10^{-2}$ \\
D & 4,0 & $-1$ & $8{,}0\times10^{-3}$ & $1{,}10\times10^{-2}$ & $1{,}10\times10^{-2}$ \\
\bottomrule
\multicolumn{6}{@{}p{0.95\columnwidth}@{}}{\scriptsize Para $A_i=0{,}5$ y
$1{,}5$ la BER medida fue 0; para $A_i=2{,}2$, $2{,}5$, 3 y 5 fue
$1{,}1\times10^{-2}$.}
\end{tabular}
\end{table}

\begin{figure}[!t]
\centering
\begin{tikzpicture}
\begin{axis}[
  width=\columnwidth, height=0.64\columnwidth,
  xlabel={Amplitud del impulso $A_i$},
  ylabel={BER $\times 10^{3}$},
  xmin=0, xmax=5, ymin=-0.5, ymax=12.5,
  grid=both, grid style={line width=0.1pt, draw=gray!30},
  legend style={font=\scriptsize, at={(0.98,0.03)}, anchor=south east},
  legend cell align=left,
  tick label style={font=\scriptsize}, label style={font=\small},
  /pgf/number format/use comma,
]
\addplot[gray, thick, dashed, no marks] coordinates {(0.00,0.0000) (0.10,0.0000) (0.20,0.0000) (0.30,0.0000) (0.40,0.0000) (0.50,0.0000) (0.60,0.0000) (0.70,0.0000) (0.80,0.0000) (0.90,0.0000) (1.00,0.0000) (1.10,0.0000) (1.20,0.0000) (1.30,0.0000) (1.40,0.0000) (1.50,0.0000) (1.60,0.0000) (1.70,0.0000) (1.80,0.0003) (1.81,0.0006) (1.82,0.0013) (1.83,0.0027) (1.84,0.0055) (1.85,0.0108) (1.86,0.0204) (1.87,0.0373) (1.88,0.0656) (1.89,0.1112) (1.90,0.1820) (1.91,0.2874) (1.92,0.4384) (1.93,0.6461) (1.94,0.9206) (1.95,1.2692) (1.96,1.6948) (1.97,2.1940) (1.98,2.7566) (1.99,3.3659) (2.00,4.0000) (2.01,4.6341) (2.02,5.2434) (2.03,5.8060) (2.04,6.3052) (2.05,6.7308) (2.06,7.0794) (2.07,7.3539) (2.08,7.5616) (2.09,7.7126) (2.10,7.8180) (2.11,7.8888) (2.12,7.9344) (2.13,7.9627) (2.14,7.9796) (2.15,7.9892) (2.16,7.9945) (2.17,7.9973) (2.18,7.9987) (2.19,7.9994) (2.20,7.9997) (2.30,8.0000) (2.40,8.0000) (2.50,8.0000) (2.60,8.0000) (2.70,8.0000) (2.80,8.0000) (2.90,8.0000) (3.00,8.0000) (3.10,8.0000) (3.20,8.0000) (3.30,8.0000) (3.40,8.0000) (3.50,8.0000) (3.60,8.0000) (3.70,8.0000) (3.80,8.0000) (3.90,8.0000) (4.00,8.0000) (4.10,8.0000) (4.20,8.0000) (4.30,8.0000) (4.40,8.0000) (4.50,8.0000) (4.60,8.0000) (4.70,8.0000) (4.80,8.0000) (4.90,8.0000) (5.00,8.0000)};
\addlegendentry{Teórica, $p_-=1/2$}
\addplot[blue, thick, no marks] coordinates {(0.00,0.0000) (0.10,0.0000) (0.20,0.0000) (0.30,0.0000) (0.40,0.0000) (0.50,0.0000) (0.60,0.0000) (0.70,0.0000) (0.80,0.0000) (0.90,0.0000) (1.00,0.0000) (1.10,0.0000) (1.20,0.0000) (1.30,0.0000) (1.40,0.0000) (1.50,0.0000) (1.60,0.0000) (1.70,0.0000) (1.80,0.0003) (1.81,0.0008) (1.82,0.0018) (1.83,0.0037) (1.84,0.0076) (1.85,0.0148) (1.86,0.0281) (1.87,0.0513) (1.88,0.0902) (1.89,0.1529) (1.90,0.2503) (1.91,0.3952) (1.92,0.6028) (1.93,0.8883) (1.94,1.2658) (1.95,1.7452) (1.96,2.3304) (1.97,3.0168) (1.98,3.7904) (1.99,4.6281) (2.00,5.5000) (2.01,6.3719) (2.02,7.2096) (2.03,7.9832) (2.04,8.6696) (2.05,9.2548) (2.06,9.7342) (2.07,10.1117) (2.08,10.3972) (2.09,10.6048) (2.10,10.7497) (2.11,10.8471) (2.12,10.9098) (2.13,10.9487) (2.14,10.9719) (2.15,10.9852) (2.16,10.9924) (2.17,10.9963) (2.18,10.9982) (2.19,10.9992) (2.20,10.9997) (2.30,11.0000) (2.40,11.0000) (2.50,11.0000) (2.60,11.0000) (2.70,11.0000) (2.80,11.0000) (2.90,11.0000) (3.00,11.0000) (3.10,11.0000) (3.20,11.0000) (3.30,11.0000) (3.40,11.0000) (3.50,11.0000) (3.60,11.0000) (3.70,11.0000) (3.80,11.0000) (3.90,11.0000) (4.00,11.0000) (4.10,11.0000) (4.20,11.0000) (4.30,11.0000) (4.40,11.0000) (4.50,11.0000) (4.60,11.0000) (4.70,11.0000) (4.80,11.0000) (4.90,11.0000) (5.00,11.0000)};
\addlegendentry{Teórica, $p_-=11/16$}
\addplot[red, only marks, mark=*, mark size=1.8pt] coordinates {
  (0,0) (0.5,0) (1,0) (1.5,0) (1.8,0) (1.9,0.28) (2,5.82) (2.1,10.7)
  (2.2,11) (2.5,11) (3,11) (4,11) (5,11)
};
\addlegendentry{Medida}
\end{axis}
\end{tikzpicture}
\caption{Síntesis del comportamiento: BER frente a la amplitud del impulso. El
umbral teórico está en $A_i=A\,\mathrm{SPS}/W=2$.}
\label{fig:ber}
\end{figure}

La BER fue nula hasta $A_i=1{,}8$, pasó a $5{,}82\times10^{-3}$ en 2 y se
estabilizó en $1{,}1\times10^{-2}$ desde $2{,}2$ (Tabla~\ref{tab:ber},
Fig.~\ref{fig:ber}). La transición fue estrecha porque $\sigma_z=0{,}025$ es
pequeño, y por encima del umbral la BER no creció con $A_i$. Los valores
medidos superaron en un factor de 1,4 los previstos con bits equiprobables por
la secuencia periódica de 1000 bits: como 1000 es múltiplo de 62,5 símbolos,
los impulsos alcanzan siempre los mismos 16 símbolos, y el piso
$11/1000$ indica que 11 de ellos valen $-1$ ($p_-=11/16$). Con ese valor,
\eqref{eq:pb_imp} reproduce todas las mediciones; con datos no periódicos el
piso tendería a $\rho/2=8\times10^{-3}$.

\subsection{Contraste con la hipótesis}\label{subsec:contraste}

La hipótesis se confirmó en lo esencial: impulsos visibles desde B sin errores,
la mitad de los $-1$ alcanzados errados en C y un piso de BER en D que no creció
al duplicar la amplitud. Los valores de C y D fueron mayores que los previstos,
pero el modelo los explica con la fracción real $p_-$. La predicción sobre la
PSD no se confirmó. Para el escenario industrial, un enlace puede mostrar un
espectro y una potencia media normales y, aun así, presentar errores cada vez
que una conmutación supera la amplitud efectiva del símbolo.

%============================================================================
\section{Prueba exploratoria en hardware SDR}\label{sec:hardware}
%============================================================================

\begin{figure*}[!t]
\centering
\setlength{\tabcolsep}{2pt}
\begin{tabular}{ccc}
\captura[0.13\textwidth]{hackrf.jpg} &
\captura[0.42\textwidth]{hw_lejos.png} &
\captura[0.42\textwidth]{hw_cerca.png} \\
\scriptsize (a) Placa HackRF One &
\scriptsize (b) Placas distanciadas &
\scriptsize (c) Placas cercanas \\
\end{tabular}
\caption{Prueba de laboratorio con 16QAM entre dos HackRF One conectadas a
computadores diferentes. En (b) y (c), de arriba abajo: espectro recibido
después del desplazamiento digital (dB frente a frecuencia en MHz, centrado en
\SI{960}{\mega\hertz}), constelación recuperada tras la sincronización de
portadora y reloj (I frente a Q) y diagrama de ojo de las componentes I y Q.}
\label{fig:hw}
\end{figure*}

\subsection{Plataforma de hardware: HackRF One}\label{subsec:hackrf}

La HackRF One es un periférico de radio definida por software (SDR) de hardware
abierto, desarrollado por Great Scott Gadgets, que transmite o recibe señales
entre \SI{1}{\mega\hertz} y \SI{6}{\giga\hertz} y se conecta al computador por
USB~\cite{hackrfdocs}. La Tabla~\ref{tab:hackrf} resume sus características
principales. En recepción, la señal pasa por etapas de ganancia de RF, FI y
banda base; el sintetizador y mezclador RFFC5072 la convierte en frecuencia, el
transceptor MAX2837/MAX2839 la filtra y amplifica, el convertidor MAX5864 la
digitaliza y el microcontrolador LPC4320 envía las muestras I/Q al computador;
en transmisión el proceso se invierte~\cite{pysdr}. En la placa utilizada
(Fig.~\ref{fig:hw}a) se distinguen la sección de RF bajo el blindaje metálico,
el microcontrolador NXP, el conector SMA de antena y los conectores de reloj de
entrada y salida.

\begin{table}[!t]
\centering
\caption{Características técnicas de la HackRF One~\cite{hackrfdocs,pysdr}}
\label{tab:hackrf}
\footnotesize
\setlength{\tabcolsep}{3pt}
\begin{tabular}{@{}p{0.36\columnwidth}p{0.6\columnwidth}@{}}
\toprule
\textbf{Característica} & \textbf{Valor} \\
\midrule
Rango de frecuencia & \SI{1}{\mega\hertz} a \SI{6}{\giga\hertz} \\
Modo de operación & Semidúplex (transmite o recibe, no ambos a la vez) \\
Tasa de muestreo & 2 a 20~MS/s en cuadratura (I/Q) \\
Resolución & 8 bits por componente (ADC/DAC) \\
Potencia máx. de TX & 5 a \SI{15}{dBm} entre \SI{10}{\mega\hertz} y \SI{2170}{\mega\hertz}; menor por encima \\
Potencia máx. de RX & \SI{-5}{dBm} (un nivel mayor puede dañarla) \\
Ganancias de RX & Amp. RF 0/\SI{11}{\decibel}; LNA 0--\SI{40}{\decibel}; VGA 0--\SI{62}{\decibel} \\
Ganancias de TX & Amp. RF 0/\SI{11}{\decibel}; FI 0--\SI{47}{\decibel} \\
Reloj & Oscilador interno de \SI{25}{\mega\hertz} o referencia externa de \SI{10}{\mega\hertz} \\
Interfaz y alimentación & USB 2.0 de alta velocidad, alimentada por el bus \\
Antena & Conector SMA hembra de \SI{50}{\ohm} \\
\bottomrule
\end{tabular}
\end{table}

Estas características explican varias condiciones de la prueba. Al ser
semidúplex, cada placa solo puede transmitir o recibir, por lo que el enlace
requirió dos unidades, cada una con su computador. La frecuencia de
\SI{960}{\mega\hertz} está en la banda en que la potencia de transmisión
máxima es de 5 a \SI{15}{dBm}, y la resolución de 8 bits limita el rango
dinámico del receptor: con las placas distanciadas, la señal ocupó pocos
niveles del convertidor y quedó cerca del piso de ruido. Además, la tasa mínima
de 2~MS/s es mayor que la del sistema simulado
($f_s=\SI{256}{\kilo\samples\per\second}$), de modo que implementar el enlace
del escenario S1 en esta placa exigiría adaptar la tasa, por ejemplo
interpolando por 8 hasta \SI{2,048}{\mega\samples\per\second}.

\subsection{Transmisión de prueba con 16QAM}\label{subsec:prueba16qam}

Como parte de la validación experimental del curso, en el laboratorio se
transmitió un mensaje con modulación 16QAM entre dos radios definidos por
software HackRF One (Fig.~\ref{fig:hw}a), cada uno conectado a un computador
diferente, uno como transmisor y otro como receptor. Esta prueba no reproduce
el sistema BPSK del escenario S1: su propósito fue verificar un enlace físico
real y observar cómo cambia la señal recibida con las condiciones del enlace.
La recepción se analizó con el espectro centrado en \SI{960}{\mega\hertz}, la
constelación recuperada y el diagrama de ojo.

Con las placas distanciadas, la señal útil apenas sobresalió del piso de ruido
del espectro, la constelación formó una nube sin estructura y la amplitud de
las componentes I y Q fue del orden de $0{,}008$ (Fig.~\ref{fig:hw}b). Los
mensajes registrados confirmaban que la transmisión se estaba realizando, pero
la modulación no podía identificarse en la constelación. Al acercar las placas,
la señal en \SI{960}{\mega\hertz} se destacó con claridad sobre el piso de
ruido, la amplitud de I y Q aumentó a cerca de $0{,}15$ y la constelación mostró
los 16 grupos de 16QAM en una cuadrícula de $4\times4$, con niveles cercanos a
$\pm0{,}27$ y $\pm0{,}85$ en cada eje (Fig.~\ref{fig:hw}c). Una componente
estrecha en \SI{960,2}{\mega\hertz} apareció con el mismo nivel en ambas
capturas, lo que sugiere que no depende del enlace.

El resultado muestra que la relación señal a ruido en el receptor determina si
los símbolos pueden distinguirse: al reducir la distancia aumentó la potencia
recibida y la constelación se volvió reconocible. En 16QAM los puntos están más
próximos entre sí que en BPSK, por lo que esta modulación requiere una mayor
relación señal a ruido. Esto es coherente con el criterio que llevó a elegir
BPSK para el escenario S1, aunque la prueba no evaluó directamente la
modulación seleccionada.



% CONCLUSIONES (sin mencionar GitHub ni GNU Radio)
%============================================================================
\section{Conclusiones}\label{sec:conclusiones}
%============================================================================

\begin{enumerate}[label=\arabic*), leftmargin=*]
  \item El enlace BPSK con filtro acoplado toleró impulsos de casi el doble de
        la amplitud del símbolo sin errores: el umbral observado coincidió con
        $A\,\mathrm{SPS}/W=2$, y la constelación de decisión confirmó la
        amplitud efectiva $A_iW/\mathrm{SPS}$.
  \item Por encima del umbral la BER se estabilizó en $1{,}1\times10^{-2}$, sin
        crecer al duplicar la amplitud: la degradación impulsiva es de umbral y
        su piso depende de la fracción y polaridad de los símbolos alcanzados,
        no de la relación señal a ruido gaussiano.
  \item El filtro acoplado diluye los impulsos más cortos que el símbolo, pero
        hace que todo impulso influya en alguna decisión.
  \item La señal temporal, la constelación de decisión, el ojo y la BER
        revelaron la perturbación; la PSD no, por lo que las medidas promediadas
        no bastan para diagnosticar ruido impulsivo escaso.
  \item El modelo reprodujo las mediciones al considerar la secuencia de prueba
        periódica, que debe tratarse como parte del diseño experimental.
  \item La prueba en hardware con 16QAM mostró que, en un enlace físico, la
        relación señal a ruido en el receptor determina si la constelación es
        reconocible, lo que refuerza la importancia de la separación entre
        símbolos considerada al elegir la modulación.
\end{enumerate}

\textbf{Trabajo futuro.} Sería ideal implementar en hardware el sistema
completo en el que se basó este proyecto: el enlace BPSK del escenario S1 con
su receptor de filtro acoplado, transmitido entre radios HackRF One, para
analizar experimentalmente el desempeño de la modulación seleccionada. Esa
implementación permitiría medir la curva de BER frente a $E_b/N_0$ con la
instrumentación de la fase de validación del curso, introducir perturbaciones
impulsivas reales o inyectadas y compararla con 16QAM bajo las mismas
condiciones. También se propone modelar llegadas aleatorias de impulsos
(Middleton o Bernoulli-gaussiano), estudiar impulsos más largos que un símbolo
y el pulso RRC, y evaluar el recorte o borrado de muestras y la codificación con
entrelazado.


\Needspace*{30\baselineskip}
%============================================================================
\section*{Anexo: preguntas previas a la simulación}\label{sec:anexo}
%============================================================================
\noindent\begin{minipage}{\columnwidth}
\centering
\captionof{table}{Respuestas a las preguntas de la guía del proyecto}
\label{tab:preguntas}
\scriptsize
\setlength{\tabcolsep}{3pt}
\begin{tabular}{@{}p{0.38\columnwidth}p{0.58\columnwidth}@{}}
\toprule
\textbf{Pregunta} & \textbf{Respuesta} \\
\midrule
1. Situación profesional & Nodo de telemetría de una planta de proceso que envía variables de máquina a una estación de supervisión. \\
2. Fenómeno físico & Perturbaciones impulsivas breves y de gran amplitud generadas por la conmutación de motores, contactores, variadores y cargas inductivas. \\
3. Información transmitida & Secuencia binaria con variables de baja velocidad: estado, temperatura o vibración. \\
4. Modulación y razón & BPSK: símbolos antipodales con máxima separación y umbral fijo en cero; prioriza la robustez sobre la tasa. \\
5. Alternativa descartada & OOK: \SI{3}{\decibel} menos de separación para igual energía y umbral dependiente del nivel recibido. \\
6. Tasa de símbolos & $R_s=\SI{32}{\kilo\sym\per\second}$. \\
7. Tasa de bits & $R_b=R_s\log_2 2=\SI{32}{\kilo\bit\per\second}$. \\
8. Muestras por símbolo & SPS $=8$. \\
9. Tasa de muestreo & $f_s=R_s\cdot\mathrm{SPS}=\SI{256}{\kilo\samples\per\second}$. \\
10. Ancho de banda & $\approx 2R_s=\SI{64}{\kilo\hertz}$ entre primeros nulos (pulso rectangular). \\
11. Variable experimental & Amplitud del impulso $A_i$: 0; 1; 2 y 4. \\
12. Evidencias & Señal temporal, PSD, constelación sobremuestreada y de decisión, diagrama de ojo y BER. \\
13. Resultado esperado & Impulsos visibles desde el nivel suave, sin errores mientras $A_iW/\mathrm{SPS}<A$, y un piso de BER por encima de ese umbral. \\
\bottomrule
\end{tabular}
\end{minipage}

%============================================================================
\section*{Declaración de uso de herramientas de IA}
%============================================================================
Se utilizó un asistente de inteligencia artificial (Claude, de Anthropic) para
apoyar la búsqueda de referencias, la redacción, la construcción del flujograma
modificado y la elaboración de las Figs.~\ref{fig:escenario}
y~\ref{fig:bloques}. Las simulaciones fueron ejecutadas por el equipo, cuyos
integrantes revisaron y validaron parámetros, resultados y conclusiones, y
asumen la responsabilidad por el contenido.


%============================================================================
% REFERENCIAS (estilo IEEE, lista manual)
%============================================================================
\begin{thebibliography}{99}
\footnotesize

\bibitem{willig2005}
A.~Willig, K.~Matheus y A.~Wolisz, ``Wireless technology in industrial
networks,'' \emph{Proc. IEEE}, vol.~93, no.~6, pp.~1130--1151, jun. 2005,
doi: 10.1109/JPROC.2005.849717.

\bibitem{gungor2009}
V.~C. Gungor y G.~P. Hancke, ``Industrial wireless sensor networks:
Challenges, design principles, and technical approaches,'' \emph{IEEE Trans.
Ind. Electron.}, vol.~56, no.~10, pp.~4258--4265, oct. 2009,
doi: 10.1109/TIE.2009.2015754.

\bibitem{middleton1977}
D.~Middleton, ``Statistical-physical models of electromagnetic interference,''
\emph{IEEE Trans. Electromagn. Compat.}, vol.~EMC-19, no.~3, pp.~106--127,
ago. 1977, doi: 10.1109/TEMC.1977.303527.

\bibitem{blackard1993}
K.~L. Blackard, T.~S. Rappaport y C.~W. Bostian, ``Measurements and models of
radio frequency impulsive noise for indoor wireless communications,''
\emph{IEEE J. Sel. Areas Commun.}, vol.~11, no.~7, pp.~991--1001, sep. 1993,
doi: 10.1109/49.233212.

\bibitem{chilo2009}
J.~Chilo, C.~Karlsson, P.~Ängskog y P.~Stenumgaard, ``Impulsive noise
measurement methodologies for APD determination in M2M environments,'' en
\emph{Proc. IEEE Int. Symp. Electromagn. Compat. (EMC)}, Austin, TX, EE.~UU.,
2009, pp.~151--154.

\bibitem{zimmermann2002}
M.~Zimmermann y K.~Dostert, ``Analysis and modeling of impulsive noise in
broad-band powerline communications,'' \emph{IEEE Trans. Electromagn.
Compat.}, vol.~44, no.~1, pp.~249--258, feb. 2002, doi: 10.1109/15.990732.

\bibitem{ghosh1996}
M.~Ghosh, ``Analysis of the effect of impulse noise on multicarrier and single
carrier QAM systems,'' \emph{IEEE Trans. Commun.}, vol.~44, no.~2,
pp.~145--147, feb. 1996, doi: 10.1109/26.486604.

\bibitem{pighi2009}
R.~Pighi, M.~Franceschini, G.~Ferrari y R.~Raheli, ``Fundamental performance
limits of communications systems impaired by impulse noise,'' \emph{IEEE
Trans. Commun.}, vol.~57, no.~1, pp.~171--182, ene. 2009,
doi: 10.1109/TCOMM.2009.0901.060440.

\bibitem{proakis2008}
J.~G. Proakis y M.~Salehi, \emph{Digital Communications}, 5.\textsuperscript{a}
ed. Nueva York, NY, EE.~UU.: McGraw-Hill, 2008.

\bibitem{sklar2001}
B.~Sklar, \emph{Digital Communications: Fundamentals and Applications},
2.\textsuperscript{a} ed. Upper Saddle River, NJ, EE.~UU.: Prentice Hall, 2001.

\bibitem{gnuradio}
GNU Radio Project, ``GNU Radio documentation,'' [En línea]. Disponible en:
\url{https://wiki.gnuradio.org}. [Accedido: nov. 2026].


\bibitem{hackrfdocs}
Great Scott Gadgets, ``HackRF One,'' en \emph{HackRF Documentation}. [En línea].
Disponible en: \url{https://hackrf.readthedocs.io/en/latest/hackrf_one.html}.
[Accedido: nov. 2026].

\bibitem{pysdr}
M. Lichtman, ``HackRF One in Python,'' en \emph{PySDR: A Guide to SDR and DSP
using Python}. [En línea]. Disponible en:
\url{https://pysdr.org/content/hackrf.html}. [Accedido: nov. 2026].

\end{thebibliography}


\end{document}
